%---
% title: "Gửi Tôi Của Năm Mười Sáu Tuổi: Giá Như Khi Ấy Đại Số Tuyến Tính Được Viết Thế Này. Kì 5 - Ma trận sơ cấp và cách tìm nghịch đảo"
% date: 2026-09-27
% tags: [dai-so-tuyen-tinh, thpt, ma-tran, nghich-dao, khu-gauss]
% summary: "Kì 5 chuỗi đại số tuyến tính cho THPT: một phép biến đổi hàng là một phép nhân ở bên trái, cả dãi biến đổi gộp thành một phép nhân, bốn điều kiện tương đương quyết định khả nghịch, và thuật toán tìm nghịch đảo bằng cách ghép ma trận không vào bên phải."
%---
\documentclass[12pt,a4paper]{article}

% Khối khai báo dưới đây CHỈ dùng khi biên dịch PDF trực tiếp
% (Overleaf, pdfLaTeX). Blog không cần nó: Pandoc tự bỏ qua phần này.
\usepackage[utf8]{inputenc}
\usepackage[T5]{fontenc}
\usepackage[vietnamese]{babel}
\usepackage{amsmath,amssymb}
\usepackage[a4paper,margin=2.2cm]{geometry}

\begin{document}

Kì trước ta đã có nghịch đảo, và chỉ còn một chỗ chưa giải được. Ta có định nghĩa, ta có công thức cho ma trận cỡ hai, nhưng với một ma trận lớn thì làm sao biết nó có khả nghịch hay không, và nếu có thì tính ra bằng cách nào?

Câu trả lời nằm ở thuật toán khử mà ta đã dùng suốt ba kì. Nhưng để dùng nó, ta phải biết trước một điều chưa từng nói ở đâu. Phép biến đổi hàng cũng là một phép nhân.

\subsection{Một phép biến đổi hàng là một phép nhân}

Ba phép biến đổi hàng sơ cấp ta vẫn dùng từ kì một.

\begin{enumerate}
  \item nhân một hàng với một số khác \(0\);
  \item đổi chỗ hai hàng;
  \item cộng một bội của hàng này vào hàng khác.
\end{enumerate}

Xét phép thứ ba, cộng \(3\) lần hàng thứ nhất vào hàng thứ ba của một ma trận \(A\) bất kỳ cỡ \(3\times n\). Hãy hỏi hàng thứ ba mới gồm những gì. Nó là ba lần hàng thứ nhất cộng với hàng thứ ba cũ. Tức là hàng thứ ba mới là một tổ hợp tuyến tính của các hàng của \(A\), với hệ số cứng là \((3,0,1)\), và đúng ba hệ số ấy không phụ thuộc vào \(A\) là gì.

Giờ áp luật hàng--cột của kì ba vào ma trận
$$
E=\left[ \begin{array}{rrr} 1 & 0 & 0\\ 0 & 1 & 0\\ 3 & 0 & 1 \end{array} \right].
$$
Hàng thứ ba của \(EA\) là tổ hợp tuyến tính của các hàng của \(A\) với hệ số lấy từ hàng thứ ba của \(E\), tức là \(3\) lần hàng thứ nhất cộng hàng thứ ba. Hai hàng trên giữ nguyên, vì hai hàng đầu của \(E\) chỉ có một số \(1\).

Vậy \(EA\) đúng là ma trận thu được từ \(A\) bằng phép cộng \(3\) lần hàng thứ nhất vào hàng thứ ba.

Điều đáng chú ý nhất nằm ở chỗ khác. Ma trận \(E\) không đổi, và ta dùng được nó cho mọi \(A\). Không cần biết \(A\) là gì trước.

\textbf{Định nghĩa 1.} Ma trận thu được bằng cách thực hiện một phép biến đổi hàng sơ cấp trên một ma trận không được gọi là \textbf{ma trận sơ cấp} (\emph{elementary matrix}).

Ba phép biến đổi cho ba dạng ma trận sơ cấp khác nhau, theo ba cách lấy hàng thứ \(i\) của ma trận không rồi sửa nó:
$$
\left[ \begin{array}{rrr} 1 & 0 & 0\\ 0 & -3 & 0\\ 0 & 0 & 1 \end{array} \right],\qquad
\left[ \begin{array}{rrrr} 1 & 0 & 0 & 0\\ 0 & 0 & 0 & 1\\ 0 & 0 & 1 & 0\\ 0 & 1 & 0 & 0 \end{array} \right],\qquad
\left[ \begin{array}{rrr} 1 & 0 & 3\\ 0 & 1 & 0\\ 0 & 0 & 1 \end{array} \right].
$$

\textbf{Định lý 1.} Nếu \(E\) là ma trận sơ cấp thu được từ một phép biến đổi hàng trên \(I_m\), và \(A\) là ma trận cỡ \(m\times n\) bất kỳ, thì \(EA\) là ma trận thu được từ \(A\) bằng đúng phép biến đổi hàng ấy.

Cách chứng minh vừa làm ở trên áp dụng cho cả ba loại: hàng thứ \(i\) của \(E\) chỉ khác hàng thứ \(i\) của \(I_m\) ở chỗ cần thiết, nên phép biến đổi rơi đúng vào hàng thứ \(i\) và đúng như thế.

\subsection{Cả một dãi biến đổi gộp thành một phép nhân}

Một phép biến đổi hàng là một phép nhân. Vậy một dãi k phép biến đổi cũng chỉ là một phép nhân, vì ta dồn các phép nhân liên tiếp:
$$
E_k\left(\cdots\left(E_2\left(E_1A\right)\right)\cdots\right)=E_k\cdots E_2E_1A.
$$

Đặc biệt, nếu ta khử được \(A\) về ma trận không bằng dãi phép biến đổi \(E_1,E_2,\ldots,E_k\) theo thứ tự ấy, thì
$$
E_k\cdots E_2E_1A=I_n.
$$
Từ đây, dùng định lý \((AB)^{-1}=B^{-1}A^{-1}\) của kì trước:
$$
A=\left(E_k\cdots E_2E_1\right)^{-1}=E_1^{-1}E_2^{-1}\cdots E_k^{-1}. \tag{1}
$$

Vế phải của (1) là một tích của các ma trận sơ cấp, vì mỗi \(E_i\) là ma trận sơ cấp và ma trận sơ cấp thì khả nghịch. Điều đó không khó: phép biến đổi hàng nào sinh ra \(E\) thì phép biến đổi ngược lại sẽ đưa \(E\) về \(I_m\), và ta sẽ thấy ngay ở dưới.

\textbf{Mệnh đề.} Hai ma trận \(A\) và \(B\) được gọi là \textbf{tương đương hàng} (\emph{row equivalent}) nếu ta có thể đi từ ma trận này sang ma trận kia bằng một dãy phép biến đổi hàng sơ cấp.

Vì mỗi phép biến đổi hàng sơ cấp đều có một phép biến đổi ngược, nên hai ma trận tương đương hàng theo chiều này thì cũng tương đương theo chiều kia.

Còn đây là danh sách phép biến đổi hàng và phép biến đổi đi ngược lại.
\begin{center}
\begin{tabular}{ll}
Phép biến đổi trên \(I\) & Phép biến đổi đưa ma trận ấy về \(I\)\\ \hline
Nhân hàng \(i\) với \(c\neq0\) & Nhân hàng \(i\) với \(1/c\)\\
Đổi chỗ hàng \(i\) và hàng \(j\) & Đổi chỗ hàng \(i\) và hàng \(j\)\\
Cộng \(c\) lần hàng \(i\) vào hàng \(j\) & Cộng \(-c\) lần hàng \(i\) vào hàng \(j\)
\end{tabular}
\end{center}

Đọc cột phải thấy ngay ma trận sơ cấp nào cũng khả nghịch, và nghịch đảo của nó cũng là ma trận sơ cấp: đó chính là phép biến đổi ở cột phải.

\subsection{Nghịch đảo chính là dãi phép biến đổi ấy}

Từ (1), lấy nghịch đảo hai vế:
$$
A^{-1}=E_k\cdots E_2E_1. \tag{2}
$$

Đọc (2) từ phải sang trái: \(A^{-1}\) thu được bằng cách làm trên ma trận không \textbf{đúng dãi phép biến đổi ấy, theo đúng thứ tự ấy}.

Ta đã có một dãi phép biến đổi hàng đưa \(A\) về \(I_n\). Chạy y hệt dãi đó trên \(I_n\) thì ra \(A^{-1}\). Không cần làm lần thứ hai, và không cần tìm ma trận sơ cấp nào cả.

\subsection{Bốn điều kiện tương đương}

Bây giờ ta có thể trả lời câu hỏi mở đầu bài. Làm sao biết một ma trận vuông có khả nghịch hay không.

\textbf{Định lý 2.} Với \(A\) là ma trận vuông cỡ \(n\times n\), bốn điều sau tương đương nhau, tức cùng đúng hoặc cùng sai.
\begin{enumerate}
  \item[(a)] \(A\) khả nghịch.
  \item[(b)] Hệ \(A\mathbf{x}=0\) chỉ có nghiệm tầm thường.
  \item[(c)] Dạng bậc thang rút gọn của \(A\) là \(I_n\).
  \item[(d)] \(A\) là tích của các ma trận sơ cấp.
\end{enumerate}

Ta chứng minh bằng cách đi vòng một vòng. Mỗi mũi tên là một kết quả, và tới mũi tên cuối cùng thì vòng khép kín.

\blogfig{\includegraphics{/images/nghich-dao-ma-tran/dong-tuong-duong.svg}}

\emph{(a) kéo theo (b).} Giả sử \(A\) khả nghịch và \(\mathbf{x}_0\) là một nghiệm của \(A\mathbf{x}=0\). Nhân hai vế với \(A^{-1}\):
$$
A^{-1}(A\mathbf{x}_0)=A^{-1}\mathbf{0},
$$
tức là \((A^{-1}A)\mathbf{x}_0=\mathbf{0}\), tức là \(I_n\mathbf{x}_0=\mathbf{0}\), tức là \(\mathbf{x}_0=0\).

\emph{(b) kéo theo (c).} Giả sử hệ \(A\mathbf{x}=0\) chỉ có nghiệm tầm thường. Khử hệ thuần nhất này bằng thuật toán Gauss--Jordan. Cột vế phải toàn số \(0\) ngay từ đầu, nên các phép biến đổi hàng không làm nó thay đổi, và ta không cần mang cột ấy theo. Hệ đã khử có duy nhất nghiệm tầm thường, nên nó phải đọc ra là
$$
x_1=0,\quad x_2=0,\quad\ldots,\quad x_n=0,
$$
tức dạng rút gọn của \(A\) chính là \(I_n\).

\emph{(c) kéo theo (d).} Giả sử khử \(A\) về \(I_n\) được, bằng dãi phép biến đổi \(E_1,\ldots,E_k\). Theo (1), \(A=E_1^{-1}E_2^{-1}\cdots E_k^{-1}\), và đây là một tích của các ma trận sơ cấp.

\emph{(d) kéo theo (a).} Giả sử \(A\) là tích của các ma trận sơ cấp. Mỗi ma trận sơ cấp đều khả nghịch, mà tích của các ma trận khả nghịch thì khả nghịch, theo định lý 2 của kì trước.

Vòng đã khép. Từ (a) suy ra (b) suy ra (c) suy ra (d) suy ra (a), nên bốn điều kiện này tương đương nhau: chỉ cần một cái đúng thì ba cái kia cũng đúng.

Trong bốn điều kiện ấy, (c) là cái tiện nhất khi ta cần quyết định xem một ma trận có khả nghịch không, vì ta đã có thuật toán khử từ kì hai. Và (b) cho thêm một điều nữa. Khi \(A\) không khả nghịch thì hệ \(A\mathbf{x}=0\) có nghiệm không tầm thường, mà những nghiệm ấy lập thành một họ vô số nghiệm.

\subsection{Thuật toán tìm nghịch đảo}

Từ (2), ta có cách làm. Muốn \(A^{-1}\), ta tìm một dãy phép biến đổi hàng đưa \(A\) về \(I_n\), rồi chạy y hệt dãi đó trên \(I_n\).

Hai việc ấy làm cùng lúc được. Ta ghép ma trận không vào bên phải của \(A\), rồi biến đổi hàng trên cả hai nửa:
$$
[\,A\mid I_n\,]\;\longrightarrow\;[\,I_n\mid A^{-1}\,].
$$
Vì mỗi phép biến đổi hàng áp dụng cho cả hai nửa, nên nửa phải luôn là dãi phép biến đổi ấy áp dụng vào \(I_n\).

\textbf{Ví dụ 1.} Tìm nghịch đảo của
$$
A=\left[ \begin{array}{rrr} 1 & 2 & 3\\ 2 & 5 & 3\\ 1 & 0 & 8 \end{array} \right].
$$
Bắt đầu từ \([\,A\mid I_3\,]\). Lấy hàng hai trừ hai lần hàng một, và hàng ba trừ hàng một:
$$
\left[ \begin{array}{rrr|rrr} 1 & 2 & 3 & 1 & 0 & 0\\ 0 & 1 & -3 & -2 & 1 & 0\\ 0 & -2 & 5 & -1 & 0 & 1 \end{array} \right].
$$
Lấy hàng ba cộng hai lần hàng hai, rồi nhân hàng ba với \(-1\):
$$
\left[ \begin{array}{rrr|rrr} 1 & 2 & 3 & 1 & 0 & 0\\ 0 & 1 & -3 & -2 & 1 & 0\\ 0 & 0 & 1 & 5 & -2 & -1 \end{array} \right].
$$
Xóa \(z\) ở trên: lấy hàng hai cộng ba lần hàng ba, lấy hàng một trừ ba lần hàng ba:
$$
\left[ \begin{array}{rrr|rrr} 1 & 2 & 0 & -14 & 6 & 3\\ 0 & 1 & 0 & 13 & -5 & -3\\ 0 & 0 & 1 & 5 & -2 & -1 \end{array} \right].
$$
Cuối cùng lấy hàng một trừ hai lần hàng hai:
$$
\left[ \begin{array}{rrr|rrr} 1 & 0 & 0 & -40 & 16 & 9\\ 0 & 1 & 0 & 13 & -5 & -3\\ 0 & 0 & 1 & 5 & -2 & -1 \end{array} \right].
$$
Nửa trái đã là ma trận không, nên
$$
A^{-1}=\left[ \begin{array}{rrr} -40 & 16 & 9\\ 13 & -5 & -3\\ 5 & -2 & -1 \end{array} \right].
$$
Kiểm lại bằng cách nhân:
$$
AA^{-1}=\left[ \begin{array}{rrr} 1 & 2 & 3\\ 2 & 5 & 3\\ 1 & 0 & 8 \end{array} \right]\left[ \begin{array}{rrr} -40 & 16 & 9\\ 13 & -5 & -3\\ 5 & -2 & -1 \end{array} \right]=\left[ \begin{array}{rrr} 1 & 0 & 0\\ 0 & 1 & 0\\ 0 & 0 & 1 \end{array} \right]=I.
$$

\subsection{Khi nào không tìm được}

Thuật toán trên không phải lúc nào cũng chạy tới cuối. Nếu \(A\) suy biến thì không có cách nào đưa nó về \(I_n\), và ta sẽ thấy điều đó bằng một hàng toàn số \(0\) xuất hiện ở nửa trái.

\textbf{Ví dụ 2.} Xét
$$
A=\left[ \begin{array}{rrr} 1 & 2 & 3\\ 4 & 5 & 6\\ 2 & 3 & 4 \end{array} \right].
$$
Từ \([\,A\mid I_3\,]\), lấy hàng hai trừ bốn lần hàng một, và hàng ba trừ hai lần hàng một:
$$
\left[ \begin{array}{rrr|rrr} 1 & 2 & 3 & 1 & 0 & 0\\ 0 & -3 & -6 & -4 & 1 & 0\\ 0 & -1 & -2 & -2 & 0 & 1 \end{array} \right].
$$
Ba lần hàng ba, rồi lấy hàng ba trừ hàng hai:
$$
\left[ \begin{array}{rrr|rrr} 1 & 2 & 3 & 1 & 0 & 0\\ 0 & -3 & -6 & -4 & 1 & 0\\ 0 & 0 & 0 & -2 & -1 & 3 \end{array} \right].
$$
Nửa trái có một hàng toàn số \(0\). Vậy dạng rút gọn của \(A\) không thể là \(I_n\), và theo điều kiện (c) của định lý 2, \(A\) suy biến.

Ta thấy được nguyên nhân ngay từ dữ kiện. Ba hàng của \(A\) không độc lập, vì
$$
4,\ 5,\ 6\;=\;3\,(2,\ 3,\ 4)\;-\;2\,(1,\ 2,\ 3),
$$
tức hàng thứ hai là ba lần hàng thứ ba trừ hai lần hàng thứ nhất. Ma trận không có nghịch đảo vì ba hàng ấy không mang ba thông tin độc lập, mà bộ ba hàng củng cả hệ luôn độc lập khi ma trận khả nghịch.

Khi hàng toàn số \(0\) xuất hiện, ta dừng lại. Có thể xóa hàng đó ở nửa trái và dừng luôn, hoặc xóa cả ở nửa phải, đều được. Việc cần làm lúc này chỉ là ghi nhận \(A\) không có nghịch đảo.

\subsection{Hệ thuần nhất}

Kì trước ta đã gặp hệ thuần nhất, và biết ngay nó luôn có nghiệm tầm thường. Câu hỏi là ngoài nghiệm ấy ra còn gì.

Định lý 2 trả lời thẳng. Với hệ vuông, hệ thuần nhất chỉ có nghiệm tầm thường khi và chỉ khi ma trận hệ số khả nghịch.

Xét hai hệ sau, cùng số ẩn, cùng số phương trình, chỉ khác hệ số.
$$
\begin{cases}
x_1+2x_2+3x_3=0\\
2x_1+5x_2+3x_3=0\\
x_1+8x_3=0
\end{cases}
\qquad\text{và}\qquad
\begin{cases}
x_1+2x_2+3x_3=0\\
4x_1+5x_2+6x_3=0\\
2x_1+3x_2+4x_3=0.
\end{cases}
$$
Ma trận hệ số của hệ thứ nhất chính là ma trận ở Ví dụ 1, ta đã biết nó khả nghịch. Nên hệ ấy chỉ có nghiệm tầm thường.

Ma trận hệ số của hệ thứ hai chính là ma trận ở Ví dụ 2, suy biến. Nên hệ ấy có nghiệm không tầm thường. Ta tìm ra một nghiệm như sau. Khử \([A\mid0]\) bằng đúng những phép biến đổi ở Ví dụ 2, cột phải lúc này toàn số \(0\) và giữ nguyên là \(0\):
$$
\left[ \begin{array}{rrr} 1 & 2 & 3\\ 0 & -3 & -6\\ 0 & 0 & 0 \end{array} \right].
$$
Hàng cuối là phương trình \(0=0\), bỏ được. Hệ còn lại là
$$
x_1+2x_2+3x_3=0,\qquad -3x_2-6x_3=0.
$$
Hàng thứ hai cho \(x_2=-2x_3\), thay vào hàng đầu thì \(x_1=x_3\). Đặt \(x_3=t\) ta được
$$
x_1=t,\quad x_2=-2t,\quad x_3=t,\qquad t\in\mathbb R.
$$
Thử \(t=1\): \((1,-2,1)\). Cả ba phương trình đều thoả, và nó khác nghiệm tầm thường.

Công việc vừa làm chỉ là đổi câu hỏi, không đổi phép tính. Cùng một dãy phép biến đổi hàng ấy, áp lên \([A\mid I_3]\) ta được nghịch đảo, áp lên \([A\mid0]\) ta được nghiệm.

\subsection{Còn thiếu gì}

Cả bài này xoay quanh ma trận vuông. Điều kiện (c) của định lý 2 nói dạng rút gọn của ma trận vuông phải là \(I_n\), và mọi thứ trên đều dựa vào điều đó.

Nhưng phần lớn hệ phương trình ta gặp không vuông. Hai phương trình ba ẩn thì sao, có khi không tồn tại nghịch đảo để tìm, và dạng rút gọn của nó cũng không thể là ma trận không vì hình dạng không khớp. Bốn phương trình hai ẩn thì sao.

Với những hệ không vuông, nghịch đảo biến mất, còn một thứ duy nhất dùng được cho mọi hình dạng. Ta đã dùng nó suốt bốn kì rồi.

Kì sau.
\end{document}
